\documentclass{article}
\usepackage{amsmath}
\usepackage{amssymb}
\usepackage{geometry}
\geometry{margin=1in}
\usepackage{titlesec}
\titleformat{\section}{\large\bfseries}{}{0em}{}[\titlerule]
\titlespacing{\section}{0pt}{1.5ex plus .5ex minus .2ex}{0.8ex}

\title{Theoretical Models of Hippocampal Information Encoding}
\author{}
\date{}

\begin{document}
\maketitle

% ---------------------------------------------------------------------------
\section*{Overview}

\paragraph{Balance parameter $\beta_i$.}
Each neuron $i$ has a balance parameter $\beta_i \in [0,1]$ drawn from a population
distribution $P_{\mathrm{balance}}$. It controls the relative contribution of the
task-driven signal versus the place-cell signal. $\beta_i = 1$ means the neuron is
fully task-driven; $\beta_i = 0$ means it is a pure place cell.

\paragraph{Model summary.}

\medskip
\begin{center}
\begin{tabular}{p{2.8cm}p{6.5cm}p{5cm}}
    \hline
    Model & High-level equation & Description \\
    \hline
    Independent &
        $r_i = \beta_i\, r^{\mathrm{task}}_i(t) + (1-\beta_i)\tilde{r}^{\mathrm{place}}_i + \varepsilon$ &
        The task signal is independent of the animal's current location and determined strictly as a function of the time since the onset of the stimuli. Each task-influenced neuron is randomly assigned 1 of 8 response functions. The place and task responses are weighted according to the balance parameter. \\[8pt]
    Place-Dependent &
        $r_i = \beta_i\, r^{\mathrm{task}}_i(t,\tilde{r}^{\mathrm{place}}_i) + (1-\beta_i)\tilde{r}^{\mathrm{place}}_i + \varepsilon$ &
        The task signal is dependent on the animal's current location and determined as a function of the time since the onset of the stimuli and the active response in the place field. Each task-influenced neuron is randomly assigned 1 of 8 response functions. The place and task responses are still weighted according to the balance parameter. \\[8pt]
    Arousal-Mediated &
        $r_i = \tilde{r}^{\mathrm{place}}_i + r^{\mathrm{task}}_i(t,\tilde{r}^{\mathrm{place}}_i) + \varepsilon$ &
        The task responses are selected from a subset of 4 functions based on whether the rat is running and whether the cell is in a place field. Each of these functions then depends solely on the time since the onset of the stimuli. \\
    \hline
\end{tabular}
\end{center}

\paragraph{Key symbols.}
\begin{itemize}
    \item $\tilde{r}^{\mathrm{place}}_i$ --- velocity-modulated Gaussian place cell rate
    \item $r^{\mathrm{task}}_i(t)$ --- TEBC task response for neuron $i$, drawn from type $k \in \{1,\dots,8\}$ (independent/place-dependent) or $k \in \{1,\dots,4\}$ (arousal-mediated)
    \item $\beta_i$ --- per-neuron balance parameter
    \item $b_i \equiv \tilde{r}^{\mathrm{place}}_i(t)$ --- live place rate, used as baseline in place-dependent and arousal-mediated models
    \item $\varepsilon \sim \mathcal{N}(-\tfrac{0.02}{30}, \tfrac{0.02}{30})$ --- additive noise term common to all models
\end{itemize}

\paragraph{Notation.} All firing rates are in Hz. Time $t$ is seconds since CS onset.
Constants: $t_{CS} = 0.25\,\text{s}$, $t_{US} = 0.75\,\text{s}$, $d_{US} = 0.1\,\text{s}$.

% ---------------------------------------------------------------------------
\section{Place Cell Firing Rate (RatInABox)}

All models inherit the Gaussian place cell implementation from the RatInABox library.
For each neuron $i$ with field centre $\mathbf{c}_i$ and width $\sigma$:

\begin{equation}
    r^{\mathrm{PC}}_i(\mathbf{x}) =
        \exp\!\left(-\frac{d(\mathbf{x},\,\mathbf{c}_i)^2}{2\sigma^2}\right)
        \cdot (f_{\max} - f_{\min}) + f_{\min}
\end{equation}

\noindent
where $d(\mathbf{x}, \mathbf{c}_i)$ is the geodesic distance from the agent position
$\mathbf{x}$ to the place field centre and $\sigma = 0.20\,\text{m}$.
The values of $f_{\min}$ and $f_{\max}$ differ by model:

\medskip
\begin{center}
\begin{tabular}{llll}
    \hline
    Model & $f_{\min}$ & $f_{\max}$ & Note \\
    \hline
    Independent      & $0\,\text{Hz}$ & $1\,\text{Hz}$  & biologically low; task constants tuned to match \\
    Place-Dependent  & $0\,\text{Hz}$ & $12\,\text{Hz}$ & physiologically realistic \\
    Arousal-Mediated & $0\,\text{Hz}$ & $12\,\text{Hz}$ & physiologically realistic \\
    \hline
\end{tabular}
\end{center}

% ---------------------------------------------------------------------------
\section{Independent Model}

\textit{The task signal is independent of the animal's current location and determined
strictly as a function of the time since the onset of the stimuli. Each task-influenced
neuron is randomly assigned 1 of 8 response functions. The place and task responses are
weighted according to the balance parameter.}

\bigskip\noindent
Each neuron $i$ has a firing rate that additively combines a velocity-modulated place cell
component and a task-driven TEBC component, with a small noise term:

\begin{equation}
    r_i(t) = \beta_i\, r^{\mathrm{task}}_i(t)
           + \bigl(1 - \beta_i\bigr)\, \tilde{r}^{\mathrm{place}}_i(t)
           + \varepsilon, \qquad
    \varepsilon \sim \mathcal{N}\!\left(-\tfrac{0.02}{30},\; \tfrac{0.02}{30}\right)
\end{equation}

\noindent
where $\beta_i \in [0, 1]$ is the per-neuron balance parameter drawn from
$\{\beta_i\} \sim P_{\mathrm{balance}}$.

\paragraph{Velocity-modulated place cell rate.}
Let $r^{\mathrm{PC}}_i$ be the standard Gaussian place cell firing rate and $v$ the
agent's smoothed speed (m/s). The modulated rate is:

\begin{equation}
    \tilde{r}^{\mathrm{place}}_i =
    \begin{cases}
        0 & v < 0.02 \\[4pt]
        \dfrac{p(100v)}{30}\; r^{\mathrm{PC}}_i & v \geq 0.02
    \end{cases}
\end{equation}

\noindent
where $p$ is the empirical cubic polynomial

\begin{equation}
    p(s) = -3.261 \times 10^{-4}\, s^3
          + 1.741 \times 10^{-2}\, s^2
          + 8.366 \times 10^{-2}\, s
          + 1.161
\end{equation}

\noindent
(argument $s = 100v$ converts m/s to cm/s).
Non-place-cell neurons have their velocity-modulated rate fixed at $\tfrac{0.02}{30}$.

\paragraph{Task firing rate.}
\begin{equation}
    r^{\mathrm{task}}_i(t) =
    \begin{cases}
        R_{\,c_i}(t) & \text{neuron } i \text{ is TEBC-responsive} \\
        0            & \text{otherwise}
    \end{cases}
\end{equation}

\noindent
where $c_i \in \{1,\dots,8\}$ is the cell type assigned to neuron $i$.
% TODO: validate the units of the type 1--8 response functions in the independent model
The response profiles $R_k(t)$ map to the equations below as follows:

\medskip
\begin{center}
\begin{tabular}{cll}
    \hline
    $k$ & Description & Equation \\
    \hline
    1 & Bimodal Gaussian          & \eqref{eq:type1} \\
    2 & One-sided Gaussian        & \eqref{eq:type2} \\
    3 & Piecewise Linear Decay    & \eqref{eq:type3} \\
    4 & Ramp with Gaussian Burst  & \eqref{eq:type4} \\
    5 & Trapezoidal US Response   & \eqref{eq:type5} \\
    6 & Noisy Uniform             & \eqref{eq:type6} \\
    7 & Windowed Bimodal Gaussian & \eqref{eq:type7} \\
    8 & Post-CS Piecewise Decay   & \eqref{eq:type8} \\
    \hline
\end{tabular}
\end{center}

% ---------------------------------------------------------------------------
\section{Place-Dependent Model}

\textit{The task signal is dependent on the animal's current location and determined as a
function of the time since the onset of the stimuli and the active response in the place
field. Each task-influenced neuron is randomly assigned 1 of 8 response functions. The
place and task responses are still weighted according to the balance parameter.}

\bigskip\noindent
The firing rate blends task and place components with the same balance weight, but the
task rate is explicitly a function of the current place rate $\tilde{r}^{\mathrm{place}}_i$:

\begin{equation}
    r_i(t) = \beta_i\, r^{\mathrm{task}}_i\!\bigl(t,\,\tilde{r}^{\mathrm{place}}_i(t)\bigr)
           + \bigl(1 - \beta_i\bigr)\, \tilde{r}^{\mathrm{place}}_i(t)
           + \varepsilon, \qquad
    \varepsilon \sim \mathcal{N}\!\left(-\tfrac{0.02}{30},\; \tfrac{0.02}{30}\right)
\end{equation}

\noindent
where $b_i \equiv \tilde{r}^{\mathrm{place}}_i(t)$ is the velocity-modulated place rate:

\begin{equation}
    r^{\mathrm{task}}_i(t) =
    \begin{cases}
        R_{c_i}\!\bigl(t,\; b_i\bigr) & \text{neuron } i \text{ is TEBC-responsive} \\
        0 & \text{otherwise}
    \end{cases}
\end{equation}

\noindent
Each profile $R_k(t, b)$ has the same functional form as in the independent model
(Equations~\eqref{eq:type1}--\eqref{eq:type8}), but all magnitude parameters are
expressed as multiples of the live baseline $b$:

\medskip
\begin{center}
\begin{tabular}{clll}
    \hline
    $k$ & Description & Magnitude parameters & Eq. \\
    \hline
    1 & Bimodal Gaussian         & $A_1 = 2.14\,b$,\quad $A_2 = 5.71\,b$                       & \eqref{eq:type1} \\
    2 & One-sided Gaussian       & $A = 3.14\,b$                                                & \eqref{eq:type2} \\
    3 & Piecewise Linear Decay   & $d_1 = 0.33\,b$,\quad $d_2 = 0.17\,b$                       & \eqref{eq:type3} \\
    4 & Ramp with Gaussian Burst & $A_{CS} = 0.40\,b$,\quad $A_{US} = 1.10\,b$                 & \eqref{eq:type4} \\
    5 & Trapezoidal US Response  & $A_{US} = 3.14\,b$                                           & \eqref{eq:type5} \\
    6 & Noisy Uniform            & (unchanged)                                                  & \eqref{eq:type6} \\
    7 & Windowed Bimodal Gaussian& $A_1 = 1.67\,b$,\quad $A_2 = 2.14\,b$                       & \eqref{eq:type7} \\
    8 & Post-CS Piecewise Decay  & steps at $0.46\,b$,\; $0.62\,b$,\; $0.38\,b$                & \eqref{eq:type8} \\
    \hline
\end{tabular}
\end{center}

\paragraph{Velocity modulation.}
The same cubic polynomial $p(s)$ is used, but below the speed threshold the rate is
scaled rather than zeroed:

\begin{equation}
    \tilde{r}^{\mathrm{place}}_i =
    \begin{cases}
        \tfrac{0.02}{30}\; r^{\mathrm{PC}}_i & v < 0.02 \\[4pt]
        \dfrac{p(100v)}{30}\; r^{\mathrm{PC}}_i & v \geq 0.02
    \end{cases}
\end{equation}

% ---------------------------------------------------------------------------
\section{TEBC Response Profiles (Types 1--8)}

The following profiles are shared by the independent and place-dependent models.
In the independent model the magnitude parameters are fixed constants (given as defaults below).
In the place-dependent model they are expressed as multiples of the live baseline
$b_i = \tilde{r}^{\mathrm{place}}_i(t)$ (see the place-dependent model table above).

% ---------------------------------------------------------------------------
\subsection*{Type 1 --- Bimodal Gaussian}

\begin{equation}\label{eq:type1}
    r(t) = A_1 \exp\!\left(-\frac{(t-\mu_1)^2}{2\sigma_1^2}\right)
         + A_2 \exp\!\left(-\frac{(t-\mu_2)^2}{2\sigma_2^2}\right) + 2b
\end{equation}

\noindent
Default parameters: $A_1 = \tfrac{0.15}{24}$, $\mu_1 = 0.3$, $\sigma_1 = 0.1$;
$A_2 = \tfrac{0.40}{24}$, $\mu_2 = 0.85$, $\sigma_2 = 0.1$; $b = \tfrac{0.07}{24}$.

% ---------------------------------------------------------------------------
\subsection*{Type 2 --- One-sided Gaussian}

\begin{equation}\label{eq:type2}
    r(t) = \begin{cases}
        A \exp\!\left(-\dfrac{(t - \mu)^2}{2\sigma^2}\right) + b & t > \mu - 0.1 \\[6pt]
        b & \text{otherwise}
    \end{cases}
\end{equation}

\noindent
Default parameters: $A = \tfrac{0.22}{24}$, $\mu = 0.75$, $\sigma = 0.15$, $b = \tfrac{0.07}{24}$.

% ---------------------------------------------------------------------------
\subsection*{Type 3 --- Piecewise Linear Decay}

Let $d_1 = \tfrac{0.02}{24}$ and $d_2 = \tfrac{0.01}{24}$.

\begin{equation}\label{eq:type3}
    r(t) = \begin{cases}
        b                                                              & t \leq 0            \\[4pt]
        \dfrac{d_1 - b}{0.2}\,t + b                                   & 0 < t < 0.2         \\[8pt]
        \dfrac{d_2 - d_1}{0.65}(t - 0.2) + 2d_1                      & 0.2 \leq t < 0.85   \\[6pt]
        d_2                                                            & 0.85 \leq t \leq 3.85 \\[4pt]
        b                                                              & t > 3.85
    \end{cases}
\end{equation}

\noindent
Default parameters: $b = \tfrac{0.06}{24}$.

% ---------------------------------------------------------------------------
\subsection*{Type 4 --- Linear Ramp with Gaussian Burst}

Let $A_{CS} = \tfrac{0.03}{24}$, $A_{US} = \tfrac{0.08}{24}$, $\sigma_{US} = 0.2$.
Define the unclamped response:

\begin{equation}
    \tilde{r}(t) = \begin{cases}
        b                                                                                     & t \leq 0.05          \\[4pt]
        \dfrac{A_{CS} - b}{0.73}(t - 0.05) + b                                              & 0.05 < t \leq 0.78   \\[8pt]
        A_{US}\exp\!\left(-\dfrac{(t - 0.78)^2}{5\,\sigma_{US}^5}\right) + b               & 0.78 < t \leq 1      \\[6pt]
        A_{CS}                                                                               & 1 < t \leq 3         \\[4pt]
        b                                                                                     & t > 3
    \end{cases}
\end{equation}

The output is clamped to $[A_{CS},\, A_{US}]$:
\begin{equation}\label{eq:type4}
    r(t) = \min\!\bigl(A_{US},\; \max\!\bigl(A_{CS},\; \tilde{r}(t)\bigr)\bigr)
\end{equation}

\noindent
Default parameters: $b = \tfrac{0.075}{24}$.

% ---------------------------------------------------------------------------
\subsection*{Type 5 --- Trapezoidal US Response}

\begin{equation}\label{eq:type5}
    r(t) = \begin{cases}
        b                                                                    & t < 0.75            \\[4pt]
        A_{US}                                                               & 0.75 \leq t \leq 0.78 \\[4pt]
        A_{US} + \dfrac{b - A_{US}}{0.42}(t - 0.78)                        & 0.78 < t < 1.2      \\[8pt]
        b                                                                    & t \geq 1.2
    \end{cases}
\end{equation}

\noindent
Default parameters: $b = \tfrac{0.035}{24}$, $A_{US} = \tfrac{0.11}{24}$.

% ---------------------------------------------------------------------------
\subsection*{Type 6 --- Noisy Uniform}

\begin{equation}\label{eq:type6}
    r(t) = b + \varepsilon, \qquad \varepsilon \sim \mathcal{N}\!\left(-\eta,\; \eta\right), \quad \eta = \frac{0.03}{30}
\end{equation}

\noindent
Default parameters: $b = \tfrac{0.04}{24}$.

% ---------------------------------------------------------------------------
\subsection*{Type 7 --- Windowed Bimodal Gaussian}

Each Gaussian peak is truncated to a one-sided window:

\begin{equation}
    G_i(t) = \begin{cases}
        A_i \exp\!\left(-\dfrac{(t - \mu_i)^2}{2\sigma_i^2}\right) & \mu_i - 3\sigma_i \leq t \leq \mu_i + D_i \\[6pt]
        0 & \text{otherwise}
    \end{cases}
\end{equation}

The total response is:

\begin{equation}\label{eq:type7}
    r(t) = \max\!\left(\sum_{i=1}^{2} G_i(t) + b,\; b\right)
\end{equation}

\noindent
Default parameters: $A_1 = \tfrac{0.07}{24}$, $\mu_1 = 0.30$, $\sigma_1 = 0.08$, $D_1 = 0.25$;
$A_2 = \tfrac{0.09}{24}$, $\mu_2 = 0.80$, $\sigma_2 = 0.15$, $D_2 = 0.10$; $b = \tfrac{0.042}{24}$.

% ---------------------------------------------------------------------------
\subsection*{Type 8 --- Post-CS Piecewise Decay}

\begin{equation}\label{eq:type8}
    r(t) = \begin{cases}
        b                                                                        & t \leq 0.25              \\[4pt]
        \tfrac{0.03}{24}                                                         & 0.25 < t \leq 0.252      \\[6pt]
        \dfrac{\tfrac{0.04}{24} - b}{0.498}(t - 0.252) + b                     & 0.252 < t < 0.75         \\[8pt]
        \dfrac{\tfrac{0.025}{24} - b}{0.75}(t - 0.75) + b                      & 0.75 \leq t \leq 1.5     \\[6pt]
        b                                                                        & t > 1.5
    \end{cases}
\end{equation}

\noindent
Default parameters: $b = \tfrac{0.065}{24}$.

% ---------------------------------------------------------------------------
\section{Arousal-Mediated Model}

\textit{The task responses are selected from a subset of 4 functions based on whether
the rat is running and whether the cell is in a place field. Each of these functions then depends solely on the time since the onset
of the stimuli.}

\bigskip\noindent
There is no balance parameter. The task rate is added directly to the
velocity-modulated place rate:

\begin{equation}
    r_i(t) = \tilde{r}^{\mathrm{place}}_i(t) + r^{\mathrm{task}}_i(t) + \varepsilon,
    \qquad \varepsilon \sim \mathcal{N}\!\left(-\tfrac{0.02}{30},\; \tfrac{0.02}{30}\right)
\end{equation}

\paragraph{Velocity modulation.}
Uses the same cubic polynomial $p(s)$, but the slow-speed scaling constant differs:

\begin{equation}
    \tilde{r}^{\mathrm{place}}_i =
    \begin{cases}
        \tfrac{0.2}{30}\; r^{\mathrm{PC}}_i & v < 0.02 \\[4pt]
        \dfrac{p(100v)}{30}\; r^{\mathrm{PC}}_i & v \geq 0.02
    \end{cases}
\end{equation}

\paragraph{Task firing rate.}
The response function is selected by two binary conditions evaluated each timestep:
whether the neuron is inside its place field, and whether the animal is running.
Place-field membership is assessed from the \emph{raw} (pre-velocity-modulation)
place rate $r^{\mathrm{PC}}_i$:

\begin{equation}
    \text{in\_field}_i = \bigl[r^{\mathrm{PC}}_i \geq 0.2\bigr], \qquad
    \text{running} = \bigl[v > 0.02\bigr]
\end{equation}

\begin{equation}
    r^{\mathrm{task}}_i(t) =
    \begin{cases}
        0 & \text{neuron } i \text{ not TEBC-responsive} \\[4pt]
        R_1\!\bigl(t,\, b_i\bigr) & \text{in field},\quad \text{running} \\[4pt]
        R_2\!\bigl(t,\, b_i\bigr) & \text{out of field},\quad \text{running} \\[4pt]
        R_3\!\bigl(t,\, b_i\bigr) & \text{in field},\quad \text{stopped} \\[4pt]
        R_4\!\bigl(t,\, b_i\bigr) & \text{out of field},\quad \text{stopped}
    \end{cases}
\end{equation}

\noindent
where $b_i = \tilde{r}^{\mathrm{place}}_i(t)$ is the velocity-modulated place rate
used as the baseline for all magnitude expressions.

\paragraph{Response functions.}

\medskip\noindent
\textbf{$R_1$ --- In field, running:}
\begin{equation}
    R_1(t, b) = \begin{cases}
        0.9\,b  & t < t_{CS}                       \\[4pt]
        b       & t_{CS} \leq t < t_{US}            \\[4pt]
        0.5\,b  & t = t_{US}                        \\[4pt]
        1.2\,b  & t_{US} < t < t_{US} + 0.1         \\[4pt]
        b       & t \geq t_{US} + 0.1
    \end{cases}
\end{equation}

\medskip\noindent
\textbf{$R_2$ --- Out of field, running:}
\begin{equation}
    R_2(t, b) = \begin{cases}
        1.1\,b  & t < t_{US}   \\[4pt]
        b       & t \geq t_{US}
    \end{cases}
\end{equation}

\medskip\noindent
\textbf{$R_3$ --- In field, stopped:}
\begin{equation}
    R_3(t, b) = \begin{cases}
        1.5\,b                                                               & t < t_{CS}              \\[6pt]
        b\!\left(1 + 0.5\,\dfrac{t_{US} - t}{t_{US} - t_{CS}}\right)       & t_{CS} \leq t < t_{US}  \\[8pt]
        3\,b                                                                 & t = t_{US}              \\[4pt]
        b                                                                    & t > t_{US}
    \end{cases}
\end{equation}

\medskip\noindent
\textbf{$R_4$ --- Out of field, stopped:}
\begin{equation}
    R_4(t, b) = b + \varepsilon_4, \qquad \varepsilon_4 \sim \mathrm{Uniform}(-0.2\,b,\; 0.2\,b)
\end{equation}

\noindent
\emph{Note:} $R_1$ and $R_3$ include responses conditioned on $t = t_{US}$ exactly,
which will not trigger under floating-point time discretisation.

\end{document}
